\documentclass[12pt]{article}
\usepackage[brazil]{babel}
\usepackage{amsmath,amsthm,amssymb,amsfonts}
\everymath{\displaystyle}
\usepackage{sbc-template}
\usepackage{graphicx,url}
\usepackage{enumitem}
\usepackage{xcolor}
\usepackage{makecell}
\usepackage[utf8]{inputenc}
\usepackage[T1]{fontenc}
\usepackage[brazil]{babel}
\usepackage{xurl}
% \usepackage{hyperref}




\usepackage{tikz}
\usetikzlibrary{shapes.geometric, arrows}

% Definição de estilos dos blocos
\tikzstyle{startstop} = [rectangle, rounded corners, minimum width=3cm, minimum height=1cm,text centered, draw=black, fill=red!30]
\tikzstyle{process} = [rectangle, minimum width=3cm, minimum height=1cm, text centered, draw=black, fill=orange!30]
\tikzstyle{decision} = [diamond, minimum width=3cm, minimum height=1cm, text centered, draw=black, fill=green!30]
\tikzstyle{arrow} = [thick,->,>=stealth]
\tikzstyle{io} = [trapezium, trapezium left angle=70, trapezium right angle=110, minimum width=3cm, minimum height=1cm, text centered, draw=black, fill=blue!30]

\sloppy

\title{Aplicação de Visão Computacional para Prevenção\\ de Acidentes em Ambientes Escolares:\\ Uma Prova de Conceito (PoC)}

\author{Brendo A. Soares\inst{1}, Wener B. de Sampaio\inst{2}, Martony D. da Silva\inst{3} }


\address{Centro de Educação Aberta e a Distância -- Universidade Federal do Piauí
  (UFPI)\\
  Caixa Postal 64049-550 -- Teresina -- PI -- Brazil
  \email{brendo.soares@ufpi.edu.br, wener.sampaio@ufpi.edu.br, mardemes@gmail.com}
}

\begin{document} 

\maketitle

\begin{abstract}
  This work presents a Proof of Concept (PoC) aimed at detecting risk situations,
  such as falls and abrupt actions, in school environments. To overcome ethical
  and legal restrictions regarding the use of images of minors, the system was
  validated using a heterogeneous dataset comprising 79 continuous events
  evaluated under an Event-Level paradigm. The architecture combines YOLOv1
  models with a bimodal heuristic, uniting postural analysis and scalar velocity.
  The results demonstrate that the velocity filter drastically reduces false
  positives, highlighting the YOLOv11s model as the ideal performance sweet spot.
  However, depth variations (Z-axis) caused limitations in recall. Thus, pose
  estimation and scale normalization emerge as guidelines for future work,
  concluding that the approach is viable as a high-precision preventive tool.

\vspace{\baselineskip}
\noindent
\textbf{Keywords}: School, Safety, Computer Vision, YOLOv11, Risk Analysis, Event-Level Evoluation.
\end{abstract}
     
\begin{resumo}
  Este trabalho apresenta uma Prova de Conceito (PoC) para detectar situações de risco,
  como quedas e ações bruscas, em ambientes escolares. Para contornar restrições éticas
  e legais sobre imagens de menores, o sistema foi validado com um conjunto heterogêneo
  de 79 eventos contínuos, avaliados sob o paradigma Event-Level. A arquitetura combina
  modelos YOLOv11 com uma heurística bimodal, unindo análise postural e velocidade escalar.
  Os resultados mostram que o filtro de velocidade reduz drasticamente os falsos positivos,
  destacando o modelo YOLOv11s como o ponto de equilíbrio ideal. Contudo, variações de
  profundidade (eixo Z) causaram limitações de revocação (Recall). Assim, a estimativa
  de pose e a normalização de escala surgem como diretrizes para trabalhos futuros,
  concluindo que a abordagem é viável como ferramenta preventiva de alta precisão.

\vspace{\baselineskip}
\noindent
\textbf{Palavras-chave}: Escola, Segurança, Visão Computacional, YOLOv11, Análise de Risco, Avaliação por eventos.
\end{resumo}



\section{Introdução} \label{introducao}

O ambiente escolar é frequentemente percebido como um local de proteção absoluta; no entanto, 
estatísticas indicam uma realidade alarmante quanto à segurança física dos alunos. Segundo
levantamento do \cite{aldeias:24}, apenas em 2023, foram registrados mais de 110 mil internações
infantis, onde as quedas representam quase 60\% dos atendimentos de emergência. Dados complementares
\cite{agenciabrasil:25} revelam que, no Brasil, cerca de 14 crianças são hospitalizadas por hora 
devido a acidentes que poderiam ser evitados com monitoramento adequado.

Apesar da presença de inspetores e professores, a supervisão humana é limitada pela fadiga, pelo
número elevado de alunos e pela impossibilidade de monitorar simultaneamente múltiplos pontos cegos,
como corredores e pátios amplos. Somando a isso, episódios de convivência violenta e \textit{bullying}
atingem 38\% das instituições \cite{abrace:25}, frequentemente manifestando-se através de dinâmica
acelerada e imprevisível dos alunos, o que potencializa o risco de quedas. Identificar essas ações de 
risco de maneira autônoma torna-se vital para mitigar a gravidade de acidentes e permitir uma resposta
imediata das equipes de segurança escolar.

A motivação para este trabalho reside na necessidade de uma camada de segurança tecnológica que atue de
forma preditiva. Sistemas de monitoramento convencionais são majoritariamente passivos, servindo apenas
para auditoria após o ocorrido. A aplicação de Visão Computacional (VC) permite a transição para um 
monitoramento ativo, capaz de identificar padrões de risco em tempo real. Contudo, muitas abordagens de
detecção de anomalias focam em ambientes industriais ou urbanos, cujos padrões de movimento de adultos
diferem drasticamente da dinâmica escolar. Além disso, a implementação prática nesses ambientes enfrenta
desafios éticos rigorosos, visto que o uso de imagens reais de menores é restrito pela Lei Geral de
Proteção de Dados (LGPD) \cite{lgpd:18} e pelo Estatuto da Criança e do Adolescente (ECA) \cite{eca:90}.

Neste contexto, este artigo descreve a implementação de uma Prova de Conceito (\textit{Proof of Concept - PoC})
que utiliza dados simulados para contornar restrições legais e éticas. Para garantir a robustez dessa validação
e testar os limites do sistema sob diferentes resoluções e ângulos de profundidade, a pesquisa expande sua
análise para um conjunto heterogêneo de dados, incorporando também vídeos de domínio público da internet e 
o \textit{UR Fall Detection Dataset} \cite{kwolek2014human}. A principal contribuição reside na validação
da arquitetura \textit{YOLOv11} (explorando diferentes escalas, do modelo \textit{Nano} ao \textit{Extra Large}) 
integrada a uma análise cinemática. Esta pesquisa busca investigar a mitigação de falsos positivos em detecções
de queda, justificando-se pelo fato de que o excesso de alarmes errôneos gera ``fadiga de alarme'' no sistema
inviabilizando o seu uso na prática, ainda que o a redução de falsos negativos (quedas não detectadas) continue
sendo um desafio crítico para trabalhos futuros. Para atingir esse equilíbrio, o estudo introduz a análise de
velocidade escalar do centroide $(v)$ como fator de confirmação junto à Razão de Aspecto $(AR)$, que consiste
na proporção matemática entre a altura e a largura da caixa delimitadora do indivíduo detectado. O objetivo
geral é validar se essa combinação é suficiente para distinguir uma queda real de movimentos voluntários,
como sentar ou agachar, garantindo a segurança física no ambiente escolar de forma automatizada.

\section{Fundamentação Teórica} \label{fundamentacao-teorica}

Nesta secção apresentam-se os conceitos fundamentais teóricos que sustentam a \textit{PoC}, abrangendo desde 
processamento de imagem até a arquitetura de detecção de objetos utilizada e as heurísticas aplicadas
para a identificação de ações de risco.

\subsection{Visão Computacional e Redes Neurais Convolucionais}

A VC busca emular a capacidade do sistema visual humano, permitindo que máquinas interpretem e compreendam
informações a partir de dados visuais. O estado da arte em VC baseia-se nas Redes Neurais Convolucionais
(\textit{Convolutional Neural Networks - CNNs}), que são modelos de \textit{deep learning} compostos por 
camadas de filtros que extraem características espaciais de forma hierárquica \cite{ibm:25}.

Em um primeiro nível, a rede identifica padrões simples como bordas e texturas; em níveis subsequentes, essas
informações são combinadas para reconhecer formas complexas, como a silhueta de um corpo humano ou a orientação
de um membro. Essa capacidade de abstração é o que permite ao sistema identificar uma pessoa em diferentes 
posturas (em pé, sentada ou deitada) no ambiente escolar.

\subsection{Arquitetura YOLO (\textit{You Only Look Once})}

Diferente de algoritmos de detecção tradicionais que utilizam janelas deslizantes ou propostas de regiões em
múltiplas etapas, a família \textit{YOLO} trata a detecção como um problema de regressão único. Em uma única
passagem pela rede (\textit{single-pass/shot}), o modelo prevê simultaneamente as coordenadas das caixas 
delimitadoras (\textit{bounding boxes}) e as probabilidades de classe \cite{yolo:20}.

Neste trabalho, utiliza-se a versão v11 \cite{ultralytics_yolo11_pt}, explorando variações que vão da escala 
\textit{Nano (n)} à \textit{Extra Large (x)}. A escolha por essa arquitetura justifica-se pela alta taxa de
quadros por segundo (\textit{Frames Per Second - FPS}) alcançada, o que é mandatório para sistemas de segurança
que exigem resposta imediata em tempo real, permitindo que a pesquisa avalie o compromisso entre:

\begin{itemize}
    \item \textbf{Eficiência}: ideal para \textit{hardwares} escolares que possuem menor poder de processamento.
    
    \item \textbf{Acurácia}: Necessária para garantir que a delimitação da caixa seja precisa o suficiente
    para os cálculos de postura e velocidade.
    
\end{itemize}

\section{Trabalhos Relacionados e Análise de Lacunas}

Embora a literatura sobre segurança escolar via VC seja restrita por questões éticas e de privacidade, 
existem avanços consolidados em áreas correlatas que servem de base para esta pesquisa. A análise desses
trabalhos permite identificar limitações que justificam a implementação desta \textit{PoC}.

\subsection{Human Action Recognition (HAR) e Video Anomaly Detection (VAD)}

Pesquisas em \textit{Human Action Recognition} (HAR) \cite{Lu_2023_ICCV} utilizam bases de dados como 
HDM05 \cite{cg-2007-2}, NTU60 \cite{shahroudy2016ntu} e NTU120 \cite{Liu_2020} para treinar redes em 
tarefas de classificação de movimentos bruscos, como corridas e quedas. No entanto, esses modelos são
majoritariamente treinados com adultos e em cenários controlados, apresentando baixa performance na
dinâmica imprevisível de pátios escolares, onde a oclusão de corpos e a variação de escala dos alunos
são constantes.

Paralelamente, abordagens de \textit{Video Anomaly Detection} (VAD) \cite{Huang_2025_CVPR}, frequentemente
baseadas em \textit{autoencoders}, buscam identificar comportamentos que fogem à regularidade temporal do
ambiente. Contudo, esses sistemas tendem a ser computacionalmente pesados e de difícil interpretação para
alertas específicos de acidentes imediatos.

\subsection{Comparativo entre Abordagens de Detecção}

Para fundamentar a escolha do modelo proposto, o Quadro \ref{qua:comparacao_abordagens} apresenta um 
comparativo entre a abordagem Baseada em Pose, cujo conceitos foram extraídos do artigo \cite{delic2025}
e as duas modelagens desenvolvidas nesta pesquisa: a Estatística (Postural), que isola a análise geométrica,
e a Proposta (Bimodal), que integra os fatores geométricos e cinemáticos.

\begin{table}[ht]
    \renewcommand{\tablename}{Quadro}
    \caption{Comparação de técnicas de detecção de quedas.}
    \centering
    \begin{tabular}{|p{3.3cm}|p{2cm}|p{3.5cm}|p{4.7cm}|}
        \hline
        \textbf{Abordagem} & \textbf{\makecell{Variáveis \\ Utilizadas}} & \textbf{Vantagem} & \textbf{Principal Lacuna} \\
        \hline
        Estática (Postural) &
        Razão de Aspecto ($AR$) &
        Baixo custo computacional &
        Alto índice de falsos positivos em ações controladas (agachar/sentar). \\
        \hline
        Baseada em Pose \cite{delic2025} &
        Pontos de articulação (\textit{Keypoints}) &
        Alta precisão na postura &
        Exige alto poder de processamento, inviável para monitoramento de múltiplos alunos em tempo real. \\
        \hline
        Proposta (Bimodal) &
        $AR$ + Velocidade ($v$) &
        Equilíbrio entre precisão e FPS; redução de falsos positivos &
        Dependência de rastreamento (\textit{tracking}) estável de IDs. \\
        \hline
    \end{tabular}
    \label{qua:comparacao_abordagens}
\end{table}

A análise comparativa revela que a integração cinemática (velocidade) atua como um filtro crítico
de confirmação. Enquanto sistemas baseados apenas em postura classificariam qualquer indivíduo
posicionado horizontalmente como uma ``queda'', a solução aqui apresentada exige a ocorrência de
um vetor de velocidade brusco para validar o alerta, preenchendo a lacuna de confiabilidade 
identificada em trabalhos anteriores.

\section{Metodologia}

Nesta secção faz-se o detalhamento dos procedimentos técnicos adotados para a execução, focando 
na integração do modelo \textit{YOLOv11} e a sistemática de coleta de dados para análise.

A arquitetura lógica do sistema, conforme ilustrada na Figura \ref{fig:metodo_bimodal}, opera em um 
fluxo bimodal. Diferente de abordagens estáticas, os metadados extraídos pelo modelo \textit{YOLOv11} alimentam
simultaneamente dois ramos de análise: a geometria postural via Razão de Aspecto ($AR$) e a dinâmica 
cinemática via velocidade escalar ($v$). O alerta de risco é condicionado à interseção lógica de ambos 
os critérios, mitigando disparos falso-positivos originados por movimentos voluntários de baixa velocidade.

\begin{figure}[ht]
    \centering
    \begin{tikzpicture}[node distance=2cm, font=\small]

    % Estilos dos blocos
    \tikzstyle{start} = [rectangle, rounded corners, minimum width=2.5cm, minimum height=0.8cm, text centered, draw=black, fill=gray!10]
    \tikzstyle{process} = [rectangle, minimum width=3cm, minimum height=0.8cm, text centered, draw=black, fill=blue!10]
    \tikzstyle{decision} = [diamond, aspect=2, minimum width=3.5cm, text centered, draw=black, fill=yellow!20]
    \tikzstyle{arrow} = [thick,->,>=stealth]

    % Nós - Coluna Central
    \node (in) [start] {Entrada: \textit{Stream} de Vídeo};
    \node (det) [process, below of=in] {Detecção e \textit{Tracking}};
    \node (ext) [process, below of=det] {Extração: $w, h, C(x,y)$};
    
    % Camada de Cálculo (Afastadas horizontalmente para evitar sobreposição)
    \node (postura) [process, below left of=ext, xshift=-1.5cm, yshift=-0.5cm] {Geometria: $AR = h/w$};
    \node (vel) [process, below right of=ext, xshift=1.5cm, yshift=-0.5cm] {Cinemática: $v = \Delta C / \Delta t$};
    
    % Decisão (Aumentado o yshift para a seta não cruzar o texto)
    \node (dec) [decision, below of=ext, yshift=-2.8cm] {($AR < 0.75$) $\land$ ($v \ge 15$)?};
    
    % Saídas
    \node (alerta) [process, below of=dec, yshift=-1cm, fill=red!20] {Gatilho: Alerta de Queda};
    \node (normal) [start, right of=dec, xshift=3.5cm] {Ação Controlada};

    % Conexões com cantos ortogonais (|- ou -|) para clareza
    \draw [arrow] (in) -- (det);
    \draw [arrow] (det) -- (ext);
    
    % Setas saindo da extração para os cálculos
    \draw [arrow] (ext) -| (postura);
    \draw [arrow] (ext) -| (vel);
    
    % Setas voltando dos cálculos para a decisão
    \draw [arrow] (postura) |- (dec);
    \draw [arrow] (vel) |- (dec);
    
    % Saída da decisão
    \draw [arrow] (dec) -- node[anchor=east] {Sim} (alerta);
    \draw [arrow] (dec) -- node[anchor=south] {Não} (normal);

    \end{tikzpicture}
    \caption{Fluxograma do método bimodal. O uso de caminhos ortogonais evita a sobreposição de elementos geométricos e cinemáticos.}    
    \label{fig:metodo_bimodal}
\end{figure}

\subsection{Ferramenta e Processamento de Dados}

O desenvolvimento foi realizado em linguagem \textit{Python}, utilizando o \textit{framework} 
oficial da \cite{ultralytics_docs} para a manipulação e inferência dos modelos de VC. O fluxo de 
processamento consiste na leitura de sequências de vídeo simulados, onde cada frame é submetido
à rede neural para a extração de metadados das detecções.

Diferente de sistemas de monitoramento convencionais, o processamento foca na extração das
coordenadas das \textit{bounding boxes} para alimentar algoritmos de análise geométrica,
sem a necessidade de armazenamento ou manipulação externa de bibliotecas de interface gráfica
complexas nesta etapa de validação.

Cabe destacar que, para a execução desta \textit{PoC}, os modelos não passaram por um processo
de treinamento específico (\textit{fine-tuning}) com imagens de ambiente escolar. Foram utilizados
os pesos pré-treinados nativos do \textit{YOLOv11} treinados na base de dados \textit{COCO}, focando
exclusivamente na detecção da classe ''pessoa`` (\textit{person}). A inteligência do sistema reside
na aplicação da heurística matemática sobre os metadados gerados por essa detecção padrão.

\subsubsection{Configuração e Escalonamento dos Modelos}

Para avaliar o desempenho do sistema sob diferentes requisitos de \textit{hardware}, foram selecionadas 
variações da arquitetura \textit{YOLOv11} (lançada como estado da arte em eficiência e acurácia). A 
escolha por esta versão justifica-se pela introdução dos módulos \textit{C3k2} e pela otimização
do \textit{Spatial Pyramid Pooling - Fast} (SPPF), que permitem uma extração de características 
mais granular em ambientes escolares densos. Foram utilizadas as seguintes variações:

\begin{itemize}

    \item \textbf{Variação \textit{Nano (v11n)}}: Configurada com menor profundidade de camadas, 
    visando testar a viabilidade em sistemas de baixa latência e recursos computacionais limitados 
    (\textit{edge computing}).
    
    \item \textbf{Variação \textit{Extra Large (v11x)}}: Utilizada para estabelecer o teto de precisão 
    e acurácia na identificação das ações de risco, servindo como padrão ouro (\textit{ground truth}
    de detecção) para a validação das heurísticas.
    
\end{itemize}


\subsubsection{Detecção de Quedas Via Heurística Geométrica}

A identificação de uma queda, classificada aqui como uma ação de risco crítica em escolas,
não depende apenas da detecção do indivíduo, mas da análise da sua geometria no espaço. O sistema
utiliza as \textit{bounding boxes} (caixas delimitadoras) geradas pelo \textit{YOLO} para calcular a $(AR)$.
\\

A lógica implementada baseia-se na variação da relação entre a largura \textit{(w)} e 
a altura \textit{(h)} da caixa. Em uma postura normal, a altura é significativamente maior que a largura ($AR < 1 $).
Complementarmente, monitora-se a velocidade de deslocamento do centroide da caixa para distinguir
uma queda de uma ação controlada, como sentar-se ou deitar-se deliberadamente.

\subsubsection{Rastreamento e Extração de Características}

Para viabilizar a análise temporal das ações de risco, a detecção isolada em quadros estáticos não é suficiente.
Por isso, implementou-se o rastreamento contínuo dos indivíduos utilizando o algoritmo de 
\textit{tracking} nativo do modelo \textit{YOLO}. Essa abordagem atribui um identificador único (\textit{ID}) a
cada indivíduo detectado, permitindo monitorar sua trajetória e evolução geométrica ao longo do tempo.

Para cada quadro processado, a rede extrai as coordenadas da \textit{bounding box} $[(x1, y1); (x2, y2)]$,
a partir das quais calcula-se o centroide $C$ do indivíduo.

\\

% $C = (C_x, C_y) = \left(\frac{x_1 + x_2}{2}, \frac{y_1 + y_2}{2} \right)$

\subsection{Heurística Combinada para Identificação de Quedas}

A identificação de uma queda real no sistema depende de uma heurística bimodal, que avalia simultaneamente a
postura estática e a dinâmica do movimento. Essa combinação foi concebida para reduzir drasticamente
a taxa de falsos positivos (por exemplo, distinguir uma pessoa deitando-se intencionalmente de uma pessoa
sofrendo uma queda abrupta).

\begin{enumerate}[label=\Alph*.]

    \item \textbf{Razão de Aspecto ($AR$)}\\
    A primeira etapa avalia a geometria da \textit{bounding box}. Define-se a largura $(w)$
    e a altura $(h)$ da detecção para calcular o $(AR)$ como mostrado na fórmula $AR = \frac{h}{w}$. 
    O estado postural do indivíduo é inferido por meio de limiares pré-definidos:\\
    Se $AR < 0,75$, a postura é classificada como ``Deitado'';\\ 
    Valores entre $0,75$ e $1,2$ indicam ``Sentado/Agachado'';\\ 
    E valores superiores indicam ``Em pé''.
    \\
    
    Ressalto que os valores limiares adotados para a classificação postural constituem como uma contribuição
    empírica dos autores. Os valores foram definidos a partir da calibração inicial realizada no \textit{Dataset}
    Simulado, onde observou-se o comportamento das dimensões das \textit{bounding boxes} de pessoas em posições
    normais (em pé) e em cenários de queda.
    \\
    \item \textbf{Análise Cinemática (Velocidade)} \\
    A segunda etapa afere a velocidade de deslocamento do centroide entre o quadro atual $(t)$ e o quadro
    imediatamente anterior $(t - 1)$ para o mesmo identificador \textit{ID} rastreado nativo do \textit{YOLO}.
    A velocidade $v$, dada em \textit{pixels} por \textit{frame}, é calculada pela distância euclidiana
    percorrida nesse intervalo. Para a interpretação do sistema, a velocidade bruta foi segmentada em três
    categorias de movimento: ``Mexendo-se Pouco'' (abaixo de 1.0 \textit{px/frame}), ``Andando'' (entre 1.0
    e 30.0 \textit{px/frame}) e ``Correndo'' (acima de 30.0 \textit{px/frame})
    \begin{equation}
        v = \sqrt{(C_{x,t} - C_{x,t-1})^2 + (C_{y,t} - C_{y,t-1})^2}
    \end{equation}
    \\
    \item \textbf{Confirmação do Evento de Risco}\\
    A confirmação de uma queda verdadeira ocorre apenas quando as duas condições críticas são satisfeitas
    simultaneamente: a transição geométrica para a postura ``Deitada'' e a ocorrência de um movimento brusco 
    no momento dessa transição. Se a velocidade for muito baixa durante a mudança de postura, o sistema 
    infere que a pessoa deitou ou agachou-se de propósito.
    
    Para parametrizar essa ação brusca, o limiar de queda foi estabelecido em $v \ge 15.0$ $px/frame$), 
    representando o ponto médio da categoria de movimento ``Andando''. Dessa forma, a validação lógica 
    obedece à seguinte condição:
    
    \begin{equation}
    Queda =
    \begin{cases}
        Verdadeiro, & \text{se } Postura = \text{``Deitado''} \land v \ge 15.0 \\
        Falso, & \text{caso contrário}
    \end{cases}
    \end{equation}

    Essa lógica atua como um filtro eficaz: a postura identificada como o corpo está, enquanto o limite de
    15.0 \textit{px/frame} garante que a transição ocorreu com velocidade suficiente para ser caracterizada
    como um acidente, e não como uma ação voluntária controlada. 
    \\
    
    A definição deste limiar de velocidade é uma proposta heurística original deste trabalho. Este valor
    foi estabelecido empiricamente durante a fase de calibração no \textit{Dataset} Simulado, identificando-se
    como o ponto crítico mínimo que separa o vetor de aceleração de um corpo em queda livre e descontrolada
    dos vetores de velocidade gerados por movimentos ou desaceleração voluntárias (como o ato de sentar-se
    ou agachar-se)
    
    
\end{enumerate}


\subsection{Processamento e Estruturação de Dados (\textit{Logging})}

A validação da \textit{PoC} requer a preservação dos resultados para posterior análise estatística. 
Para tanto, o fluxo de processamento de dados não se restringe à exibição visual (\textit{overlays} 
nos vídeos), mas implementa uma rotina estruturada de gravação de \textit{logs}.

Ao iniciar a análise de um vídeo, o sistema cria automaticamente uma estrutura de diretórios isolando os dados 
de acordo com o nome do arquivo fonte e a versão do modelo \textit{YOLO} utilizado. Todos os eventos críticos identificados
pelas heurísticas são registrados em um arquivo \textit{.csv}, que atua como um banco de dados temporal com os
seguintes campos:

\begin{itemize}

    \item \textbf{Quadro}: Número do \textit{frame}.
    
    \item \textbf{Data/Hora}: \textit{Timestamp} da ocorrência.
    
    \item \textbf{Tipo de Evento}: Classificação do alerta (ex.: QUEDA).
    
    \item \textbf{Detalhes}: Descrição textual da detecção.
    
    \item \textbf{Postura Inferida}: Estado geométrico do alvo no momento do evento.
    
    \item \textbf{Velocidade}: Taxa de deslocamento registrada durante o gatilho da ação.
    
\end{itemize}
Essa abordagem não bloqueante garante que os registros das anomalias não comprometam o tempo de
processamento de \textit{FPS} do sistema e fornece material empírico preciso para a avaliação
de acurácia apresentada nos resultados.


\subsection{Base de Dados e Cenários de Validação}

Para assegurar a qualidade da \textit{PoC} e contornar restrições éticas referentes ao uso de
imagens de menores, a validação do sistema baseou-se em um conjunto heterogêneo de dados visuais.
O experimento consistiu no processamento de sequências de vídeos categorizados como eventos 
contínuos, totalizando 79 eventos distintos analisados (entre cenários normais e ações de risco).
As fontes de dados foram divididas em três categorias:

\begin{itemize}

    \item \textbf{\textit{Dataset} Simulado}: Composto por gravações controladas em ambiente escolar
    (sala de aula), produzidas exclusivamente pelos autores deste trabalho para esta pesquisa. O
    processo envolveu a gravação com câmeras posicionadas em ângulos que simulam o monitoramento de
    uma sala, capturando a dinâmica de movimentação, interações em grupo, quedas ensaiadas e ações
    voluntarias (como agachar-se). Esta base de dados contem 4 registros em plano contínuo com 
    variação de 0:58 à 2:40 minutos. Estes dados forneceram a base para a calibração inicial da 
    heurística da Postura $(AR)$ e Velocidade $(v)$.
    \\
    \item \textbf{\textit{UR Fall Detection Dataset}}: Base pública amplamente utilizada na literatura 
    acadêmica para o \textit{benchmark} de sistemas de detecção de anomalias contendo 30 eventos
    de quedas \cite{kwolek2014human}. Este \textit{dataset} é composto por gravações realizadas com
    câmeras Kinect, capturando tanto Atividade de Vida Diária (ADL) como sentar, caminhar e deitar no
    chão, quanto eventos de quedas frontais e laterias de indivíduos em posições variadas. Este
    \textit{dataset} foi fundamental para testar o comportamento do modelo frente a padrões de quedas
    já consolidados cientificamente.
    \\
    \item \textbf{Vídeos de Domínio Público}: Seleção de cenas de quedas reais extraídas da internet,
    incluindo registros de câmeras de segurança, nessa base de dados foram usados 2 eventos, 1 contendo
    um evento de queda  que foi retirado de um \textit{Reality Show} e 1 com evento normal de pessoas 
    andando. O uso deste conjunto visou introduzir ``ruídos'' do mundo real, desafiando a rede neural 
    com as variações de iluminação, qualidade de imagem e, primordialmente, mudanças de perspectiva 
    no eixo de profundidade (eixo Z).    
    
\end{itemize}

Esta composição de dados permitiu uma validação criteriosa sobre a capacidade de generalização dos modelos
da família \textit{YOLOv11}, confrontando a precisão da detecção com a diversidade de contextos de entrada.

Para melhor precisão e a uniformidade da extração de características, todos as sequências de vídeo foram
submetidas a uma compressão para a resolução nativa de 1080p, com o fito de neutralizar disparidades de 
escala espacial decorrentes de resoluções originais distintas. No tocante ao processamento, a inferência dos
\textit{datasets} foi executada em um ambiente equipado com processador Intel core i5-13450hx e 16 GB de 
memória RAM. Restringiu-se o processamento deliberadamente à Unidade Central de Processamento (\textit{CPU}),
sem aceleração gráfica dedicada. Essa premissa metodológica buscou simular as limitações de infraestrutura
inerentes à realidade da infraestrutura das escolas públicas. Embora a aferição métrica exata da taxa de
quadros por segundo (\textit{FPS}) e do consumo de \textit{hardware} não tenha composto o escopo analítico
desta \textit{PoC} que focou primordialmente na eficácia da filtragem de Falsos Positivos, a execução 
contínua puramente em \textit{CPU} demonstrou a viabilidade operacional básica do modelo bimodal.

\subsection{Estratégia de Avaliação: Análise Baseada Em Eventos}

Durante a fase de testes preliminares desta pesquisa, observou-se empiricamente que a avaliação métrica
puramente baseada em quadro (\textit{Frame-Level}) tende a penalizar severamente o desempenho estatístico
do modelo. Constatou-se que uma falha de inferência em frações de segundo não anula a eficácia do alerta já
gerado para a ação contínua. Portanto, para refletir com maior fidelidade a aplicação da \textit{PoC} em
ambientes escolares, onde o objetivo das equipes de segurança é serem notificadas sobre o incidente de
maneira total, a validação dos resultados esta pesquisa foi estruturada sob a abordagem baseada em eventos
(\textit{Event-Level}).

Com esta métrica, uma sequência contínua de ação, como o ato de uma pessoa ir ao chão é agrupada e
contabilizada como um único evento temporal. A matriz de Confusão para a extração das métricas foi
definida pelas seguintes regras:

\begin{enumerate}
    \item Verdadeiro Positivo (\textit{TP}): Um evento real de queda que disparou ao menos um alerta de 
    risco válido do sistema.
    \item Falso Negativo (\textit{FN}): Um evento real de queda que o sistema falhou em alertar durante
    toda a sua duração.
    \item Verdadeiro Negativo (\textit{TN}): Um evento normal como uma pessoa caminhando ou sentando que
    o sistema corretamente ignorou, não disparando alertas.
    \item Falso Positivo (\textit{FP}): Um evento normal que o sistema classificou de maneira errada como
    uma ação de risco.
\end{enumerate}

Essa abordagem reflete com maior fidelidade a aplicação da \textit{PoC} em ambientes escolares, onde o
objetivo das equipes de segurança é ser notificado sobre o incidente, independentemente de flutuações
de \textit{bounding boxes} em quadros intermediários.

\section{Resultado e Discussão}

Analisando estatisticamente o desempenho dos modelos \textit{YOLOv11} na identificação de ações de risco. 
Conforme definido na metodologia, a extração de métricas foi conduzida sob o padrão \textit{Event-Level}, 
contabilizando o acionamento do sistema durante a janela de ocorrência da ação contínua. O experimento 
processou eventos distintos, combinando cenários normais (Verdadeiros Negativos) e episódios de quedas reais
(Verdeiros Positivos).

\subsection{Desempenho Isolado por Cenário de Validação}

Para compreender a capacidade de generalização do sistema, os resultados foram primeiramente analisados
de forma segmentada, observando o comportamento da heurística bimodal (Postura + Velocidade) frente às
particularidades de cada base de dados:


\textbf{No Dataset Simulado (Ambiente Escolar)}: Por tratar-se de um ambiente controlado, com câmeras
posicionadas de forma a simular o monitoramento de uma sala de aula com boa resolução. Este \textit{dataset}
busca introduzir as complexidades da rotina escolar, incluindo oclusões movimentações em grupo e ações
voluntarias rotineiras, como agachar-se para pegar algum objeto. A Tabela ~\ref{tab:metricas-meu-dataset}
apresenta o comparativo das métricas neste cenário.

\setcounter{table}{0}
\begin{table}[h!]
\centering
\small
\caption{Métricas de detecção: Desempenho no \textit{Dataset} Simulado.}
\label{tab:metricas-meu-dataset}
\resizebox{\textwidth}{!}{%
\begin{tabular}{l|c|c|c|c|c|c|c|c|c}
\hline
\textbf{Cenário} &\textbf{Modelo} & \textbf{TN} & \textbf{FP} & \textbf{FN} & \textbf{TP} & \textbf{Acurácia} & \textbf{Precisão (P)} & \textbf{Recall (R)} & \textbf{F1-Score} \\ \hline
Com Velocidade & YOLOv11n & 2 & 7  & 4 & 3 & 0.312 & 0.300 & 0.428 & 0.352 \\ 
Com Velocidade & YOLOv11s & 2 & 7  & 3 & 4 & 0.375 & 0.363 & 0.571 & 0.444 \\ 
Com Velocidade & YOLOv11m & 3 & 6  & 4 & 3 & 0.375 & 0.333 & 0.428 & 0.375 \\ 
Com Velocidade & YOLOv11l & 3 & 6  & 5 & 2 & 0.312 & 0.250 & 0.285 & 0.266 \\ 
Com Velocidade & YOLOv11x & 1 & 9  & 4 & 3 & 0.235 & 0.250 & 0.428 & 0.315 \\ 
Sem Velocidade & YOLOv11m* & 2 & 16 & 2 & 5 & 0.280 & 0.238 & 0.714 & 0.357 \\ \hline
\end{tabular}%
}
\textit{*Nota: O modelo YOLOv11m representa todos os modelos da escala no cenário sem velocidade, cujos resultados foram idênticos.}
\end{table}

Neste ambiente altamente dinâmico, a ausência de análise de velocidade provou-se inviável para produção,
gerando alarmes falsos em massa (\textit{FP} = 16). A ativação da heurística reduziu drasticamente esses
erros para apenas 6 no modelo \textit{YOLOv11m}. Observa-se que o modelo \textit{YOLOv11x} apresentou a 
maior incidência de falsos positivos (\textit{FP} = 9), sugerindo que arquiteturas excessivamente profundas 
tendem a sofrer de \textit{overfitting} nas características das \textit{bounding boxes} durante oclusões, 
causando instabilidade no cálculo do centroide.

    
\textbf{No \textit{UR Fall Detection Dataset}}: Esta base pública acadêmica é caracterizada por gravações
vídeos de quedas em ambientes controlados, com excelente iluminação e com a câmera posicionada de forma a
capturar a movimentação de forma clara no eixo bidimensional. A Tabela ~\ref{tab:metricas-ur-fall} detalha
os resultados obtidos neste conjunto de dados.

\begin{table}[h!]
\centering
\small
\caption{Métricas de detecção: Desempenho no \textit{Dataset: UR Fall Detection Dataset}.}
\label{tab:metricas-ur-fall}
\resizebox{\textwidth}{!}{%
\begin{tabular}{l|c|c|c|c|c|c|c|c|c}
\hline
\textbf{Cenário} & \textbf{Modelo} & \textbf{TN} & \textbf{FP} & \textbf{FN} & \textbf{TP} & \textbf{Acurácia} & \textbf{Precisão (P)} & \textbf{Recall (R)} & \textbf{F1-Score} \\ \hline
Com Velocidade & YOLOv11n  & 30 & 0 & 18 & 12 & 0.700 & 1.000 & 0.400 & 0.571 \\
Com Velocidade & YOLOv11s  & 30 & 0 & 12 & 18 & 0.800 & 1.000 & 0.600 & 0.750 \\
Com Velocidade & YOLOv11m  & 30 & 0 & 13 & 17 & 0.783 & 1.000 & 0.566 & 0.723 \\
Com Velocidade & YOLOv11l  & 30 & 0 & 15 & 15 & 0.750 & 1.000 & 0.500 & 0.666 \\
Com Velocidade & YOLOv11x  & 30 & 0 & 15 & 15 & 0.750 & 1.000 & 0.500 & 0.666 \\
Sem Velocidade & YOLOv11m* & 28 & 2 & 10 & 20 & 0.800 & 0.909 & 0.666 & 0.769 \\ \hline
\end{tabular}%
}
\textit{*Nota: O modelo YOLOv11m representa todos os modelos da escala no cenário sem velocidade, cujos resultados foram idênticos.}
\end{table}

Ao analisar a Tabela ~\ref{tab:metricas-ur-fall} pode-se comprovar o poder do limitador cinemático em
ambientes controlados. Com a velocidade ativada, o sistema alcançou 100\% de Precisão (Zero Falsos Positivos)
em todas as variantes da rede. O modelo \textit{YOLOv11s} atingiu o melhor equilíbrio (\textit{F1-Score} de 0.750)
Embora o cenário puramente postural (Sem Velocidade) tenha um \textit{Recall} levemente maior, ele introduziu 
Falsos Positivos (\textit{FP} = 2), demonstrando que a velocidade é indispensável para a confiabilidade do alerta.


\textbf{Nos Vídeos de Domínio Público (Internet)}: Este conjunto serviu como teste de estresse focando em variação extrema de resolução e perspectiva (\textit{long shot}) e com deslocamento no eixo de profundidade (eixo Z).

\begin{table}[h!]
\centering
\small
\caption{Métricas de detecção: Desempenho no \textit{Dataset}: Vídeos Público.}
\label{tab:metricas-videos-publico}
\resizebox{\textwidth}{!}{%
\begin{tabular}{l|c|c|c|c|c|c|c|c|c}
\hline
\textbf{Cenário} & \textbf{Modelo} & \textbf{TN} & \textbf{FP} & \textbf{FN} & \textbf{TP} & \textbf{Acurácia} & \textbf{Precisão (P)} & \textbf{Recall (R)} & \textbf{F1-Score} \\ \hline
Com Velocidade & Todos & 2 & 0 & 1 & 0 & 0.666 & 0.000 & 0.000 & 0.000 \\ 
Sem Velocidade & Todos & 1 & 4 & 0 & 1 & 0.333 & 0.200 & 1.000 & 0.333 \\ \hline
\end{tabular}%
}
\end{table}

Os dados da Tabela ~\ref{tab:metricas-videos-publico} revelaram a principal vulnerabilidade do sistema cinemático
bidimensional. Com a velocidade ativada, o sistema não conseguiu classificar as quedas (\textit{Recall} nulo), pois
a distância da câmera reduziu o movimento absoluto em \textit{pixels} para um valor inferior ao limiar ($ v \ge 15.0$).
Quando a velocidade foi desligada, o \textit{YOLO} detectou perfeitamente as quedas (\textit{Recall} 100\%), provando
que a VC enxerga a ação, mas a heurística matemática a bloqueia pela falta de escala de profundidade.

\subsection{Análise Quantitativa e Comparativo de Arquiteturas}

A consolidação dos dados provenientes dos três cenários reflete a eficácia da abordagem proposta. 
A Tabela ~\ref{tab:metricas-geral} apresenta o desempenho detalhado do escalonamento da arquitetura
\textit{YOLOv11} sobre o total de eventos processados na avaliação global.

\begin{table}[ht!]
\centering
\small
\caption{Comparativo geral de métricas: Cenários com e sem limitador cinemático.}
\label{tab:metricas-geral}
\resizebox{\textwidth}{!}{%
\begin{tabular}{l|c|c|c|c|c|c|c|c|c}
\hline
\textbf{Cenário} & \textbf{Modelo} & \textbf{TN} & \textbf{FP} & \textbf{FN} & \textbf{TP} & \textbf{Acurácia} & \textbf{Precisão (P)} & \textbf{Recall (R)} & \textbf{F1-Score} \\ \hline
Com Velocidade & YOLOv11n & 34 & 7  & 23 & 15 & 0.620 & 0.681 & 0.394 & 0.500 \\ 
Com Velocidade & YOLOv11s & 34 & 7  & 16 & 22 & 0.708 & 0.758 & 0.578 & 0.656 \\ 
Com Velocidade & YOLOv11m & 34 & 6  & 18 & 20 & 0.696 & 0.769 & 0.526 & 0.625 \\ 
Com Velocidade & YOLOv11l & 35 & 6  & 18 & 17 & 0.658 & 0.739 & 0.447 & 0.557 \\ 
Com Velocidade & YOLOv11x & 33 & 9  & 20 & 18 & 0.637 & 0.666 & 0.473 & 0.553 \\ 
Sem Velocidade & YOLOv11m & 31 & 22 & 12 & 26 & 0.626 & 0.541 & 0.684 & 0.604 \\ \hline
\end{tabular}%
}
\textit{*Nota: O modelo YOLOv11m representa todos os modelos da escala no cenário sem velocidade,
cujos resultados foram idênticos.}.
\end{table}

A análise da Tabela ~\ref{tab:metricas-geral} evidencia um comportamento contra-intuitivo: o aumento da
complexidade da rede neural não se traduz em maior eficácia. O modelo \textit{YOLOv11s} antingiu a melhor 
capacidade preditiva global (\textit{F1-Score} de 0.656), enquanto o escalonamento par arquiteturas 
profundas (\textit{Large e Extra Large}) resultou em degradação das métricas. Do ponto de vista da VC, isso
sugere que redes excessivamente complecas sofrem de instabilidade (\textit{overfitting} espacial) em cenários
de oclusão comuns ao ambiente escolar. Nas arquiteturas maiores, a \textit{bounding boxes} oscila intersamente
a cada quadro tentando se ajustar aos mínimos detalhes, gerando variações artificiais na posição do centroide.
Esse ruído corrompe o cálculo cinemático de velocidade, provocando falsos positivos. Em contrapartida, modelos
mais raros e generalistas (\textit{Small} e \textit{Medium}) atuam de forma mais estável no rastreamento contínuo,
garantindo que a heurística bimodal funcione conforme o esperado.


Outro ponto importante a ser evidenciado é que, sem o limitador de velocidade, a heurística  de Razão de Aspecto
tornou-se o principal o gargalo do sistema sem velocidade. Como consequência, qualquer \textit{bounding box},
independentemente da precisão do modelo, passou a acionar o critério de detecção, tornando irrelevantes os
ganhos de precisão obtidos pelo modelo.

\subsection{Limitação Cinemáticas (O gargalo do Eixo Z)}

Embora a avaliação por eventos \textit{Event-Level} comprove a viabilidade da filtragem de ações voluntarias
(como evidenciado pelo ganho expressivo de Precisão Global), o \textit{trade-off} com a métrica de
\textit{Recall} expõe a necessidade de evolução da \textit{PoC}. O cálculo da velocidade euclidiana do
centroide ($Cx, Cy$) em valores absolutos demonstrou-se inviável para captura de movimentos no eixo de
profundidade (eixo Z)

Conclui-se que o uso de caixas delimitadoras e cálculos bidimensionais resolve o problema de falso positivo
na rotina escolar, mas futuras implementações devem buscar a normalização da velocidade de acordo com as 
dimensões da caixa ou a adoção de técnicas de \textit{Pose Estimation} para garantir a detecção espacial
correta em qualquer perspectiva.

\section{Considerações Finais}

Esta Prova de Conceito investigou a viabilidade de sistemas de Visão Computacional baseados na
arquitetura \textit{YOLOv11} para identificação de ações de risco em salas e pátios escolares,
sob uma perspectiva rigorosa de respeito à privacidade de dados. A adoção da métrica de uma
avaliação em eventos (\textit{Event-Level}) provou ser essencial na pesquisa. Ao evitar as
distorções estatísticas provocadas por flutuações momentâneas de inferência de detecção de
\textit{bounding boxes}, essa abordagem refletiu com altíssima fidelidade a aplicação prática
de um sistema de monitoramento preditivo.

Os resultados já consolidados demonstram o sucesso da abordagem bimodal (Postura + Velocidade)
na mitigação de alarmes falsos, que constitui um dos maiores gargalos operacionais na segurança
eletrônica. Em ambientes controlados e com boa resolução focal, como o \textit{UR Fall Detection Dataset},
a inserção da variável cinemática conferiu ao sistema uma precisão absoluta de 100\% de zero
falsos positivos, em todas as variantes testadas. No panorama global, que integrou a complexidade
de uma rotina escolar simulada e testes de estresse de vídeos de internet, a variante \textit{YOLOv11s}
com filtro cinemático estabeleceu o melhor ponto de equilíbrio prático, alcançando uma Acurácia
de 70,89\% e um \textit{F1-Score} de 0.656.

Por outro lado, o experimento cumpriu seu papel científico ao mapear com precisão os limites
físicos da abordagem bidimensional absoluta. A redução do \textit{Recall} observada nos vídeos
de plano aberto e de domínio público revelou que a velocidade calculada em \textit{pixels}
sem normalização de escala torna-se vulnerável a eventos que ocorrem no eixo de profundidade
(eixo Z). Essa descoberta fundamenta cientificamente a necessidade de evolução da arquitetura
do sistema para aplicações comerciais estáveis.

\subsection{Trabalhos Futuros}

Para superar os gargalos identificados e transicionar a solução para um ambiente de produção
ativa, propõem-se os seguintes desdobramentos:

\begin{enumerate}

    \item \textbf{Normalização de Escala e Profundidade}: Substituir o cálculo de velocidade em \textit{pixels}
    absolutos por uma métrica normalizada em relação à dimensão da \textit{bounding box} ou à altura média do
    indivíduo na cena. Isso tornará o sistema flexível à distância da pessoa em relação à câmera.
    \\
    \item \textbf{Integração de \textit{Pose Estimation}}: A transição da análise baseada em caixas delimitadoras 
    (\textit{bounding box}) para o rastreamento de pontos-chave do corpo (esqueletos) permitirá uma descrição
    cinemática tridimensional, eliminando a ``cegueira'' do sistema para quedas no eixo Z.
    
\end{enumerate}

Em última análise, este trabalho consolida a visão computacional como uma tecnologia viável e necessária 
para a segurança escolar. A \textit{PoC} aqui desenvolvida serve como base tecnológica e ética para que
futuras implementações possam elevar a segurança institucional, transformando a análise de dados brutos
em um ambiente de proteção ativa para a comunidade escolar.

\raggedright
\bibliographystyle{sbc}
\bibliography{sbc-template}

\end{document}
